% !TeX root = ../drafts/b-polynomial-commitment-scheme.tex
\section{The polynomial commitment scheme: WHIR \slash Ligerito}\label{annex:pcs}

A self-contained description of the WHIR\slash Ligerito polynomial commitment scheme \cite{ACFY25,NA25} over binary fields, with a round-by-round soundness analysis in the Johnson list-decoding regime. The two schemes coincide when WHIR is instantiated with interleaved codes and Ligerito with the Reed-Solomon codes of Annex~\ref{app:ntt}, which gives the novel polynomial basis and additive NTT used for encoding. Following Flock~\cite{Flock26}, we do not use out-of-domain (OOD) sampling at commitment, so our PCS is only list binding. The two fields are those of \S\ref{sec:vm}: we commit to a $\K$-valued multilinear, and opening claims are in $\E$.

\paragraph{Symbols.}

\begin{center}\small
\begin{tabular}{@{}lll@{}}
\hline
here & in \cite{ACFY25,NA25,BCHKS25} & meaning\\
\hline
$\nv_i$        & $m_i$        & variables of the multilinear at level $i$\\
$\nlev$        & $r$          & number of levels\\
$\ell_i$       & $\ell_i$, $\nu$ & fold factor, and the interleaving exponent\\
$\rate_i$      & $\rho_i$     & code rate\\
$\rexp_i$      & $R_i$        & rate exponent, $\rate_i=2^{-\rexp_i}$\\
$\rad_i$       & $\gamma_i$   & proximity radius\\
$\slack_i$     & $\eta_i$     & Johnson slack\\
$\mcapar$      & $\theta$     & MCA parameter (denoted $m$ in \cite{BCHKS25})\\
$\dmin$        & $\delta$     & relative minimum distance\\
$\lambda$      & $\lambda$, $\alpha$ & batching challenge, as in \S\ref{sec:stacking}\\
$\fc_j$        & $s_j$        & sumcheck fold challenge\\
$\rc{j}$       & $\sigma_j$   & running sumcheck claim after round $j$\\
$J_i$          & $T_i$        & number of claims, as in \S\ref{sec:e2e-unrolled}\\
$\nq_i$        & $t_i$        & query count\\
$\qset_i$      & $J_i$        & the sampled columns\\
$\blen$        & $n$          & code block length\\
$\dom$         & $S$          & evaluation domain\\
$\orc^{(i)}$   & $C^{(i)}$    & oracle, an interleaved word ($C_j$ is a constraint)\\
$\nrows$       & (none)       & rows of $\orc^{(0)}$ the prover commits, $\nrows \le 2^{\ell_0}$\\
$\nstack$      & (none)       & field elements placed in the stack, as in \S\ref{sec:stacking}\\
$\lst_{\rad}$  & $\lambda_\gamma$ & list of codewords within $\rad$\\
$W$, $c$       & $w$, $v$     & weight and value of a claim, as in \S\ref{sec:stacking}\\
$\ffinal$      & $p$          & final plaintext multilinear\\
$\cood$        & $y$          & out-of-domain answer ($y$ generates $\E$ over $\K$)\\
$\mcanum$      & $a$          & numerator of the MCA error\\
$\vansub_i$, $\Wn{i}$ & $W_i$, $\widehat W_i$ & subspace vanishing polynomial, normalized\\
$\novel_j$     & $X_j$        & novel-basis polynomial ($X_i$ is a sumcheck variable)\\
$\ntmap_d$     & $q_d$        & linearized map of the NTT ($q$ is the witness stack)\\
$\bfarr$       & $B$          & butterfly array\\
$\mpoly$       & $P$          & message polynomial ($P_i$ is a committed column)\\
$e_c$          & $\beta_c$    & the fixed $\Ftwo$-basis, as in Annex \ref{annex:rs}\\
$\npad$, $\padc_u$ & (none)   & padding size, row $u$'s padding (\S\ref{sec:pcs-zk})\\
$\padfold$, $\rlane$ & (none) & folded padding, random row (\S\ref{sec:pcs-zk})\\
\hline
\end{tabular}
\end{center}

%=====================================================================
\subsection{Preliminaries}\label{sec:prelims}
%=====================================================================

\subsubsection{Codes, interleaving, and list decoding}\label{sec:codes}

We use the Reed-Solomon code $\RS[\E,\dom,\kdim]$ of Definition~\ref{def:rs-code}, with $|\dom|=\blen$, dimension $\kdim=2^\kappa$, and rate $\rate=\kdim/\blen$. Messages are indexed by $\cube{\kappa}$; the linear encoder $\Enc$ of \S\ref{sec:binary-rs} maps them onto the code and preserves $\K$-valued words.

\begin{definition}[Interleaved words, lists, folds]\label{def:interleaved}
An \emph{interleaved word} is a matrix $\orc \in \E^{\cube{\ell} \times \dom}$; it is an \emph{interleaved codeword} if every row $\orc[u, \cdot]$ is a codeword $(u \in \cube{\ell})$. $\orc$ \emph{agrees} with an interleaved codeword $U$ on $\aset \subseteq \dom$ if $\orc[u,x] = U[u,x]$ for all $u \in \cube{\ell}$ and all $x \in \aset$; it is \emph{$\rad$-close} to $U$ if they agree on some $\aset$ with $|\aset| \ge (1-\rad)\blen$; and $\lst_\rad(\orc)$ is the set of interleaved codewords $\rad$-close to $\orc$. A single word \emph{agrees with the code} on $\aset$ if some codeword equals it there, and $\orc$ \emph{agrees with the interleaved code} on $\aset$ if some interleaved codeword equals its rows there. An interleaved codeword $U$ \emph{decodes} to the multilinear $g_U : \cube{\ell + \kappa} \to \E$ with $U[u, \cdot] = \Enc(g_U(u, \cdot))$; distinct $U$ decode to distinct $g_U$, hence distinct $\mle{g_U}$. For $\fc \in \E^{j}$, $j \le \ell$, the \emph{fold} $\orc_\fc \in \E^{\cube{\ell - j} \times \dom}$ is $\orc_\fc[u, x] := \sum_{\iota \in \cube{j}} \eq(\fc, \iota)\, \orc[(\iota, u), x]$ (for $j = \ell$, a single word in $\E^{\dom}$).
\end{definition}

\begin{fact}[Folding commutes with decoding]\label{fact:fold-decode}
If $U$ is an interleaved codeword then so is $U_\fc$, with $\mle{g_{U_\fc}} = \mle{g_U}(\fc, \cdot)$ (linearity of $\Enc$ plus Fact~\ref{fact:eq}).
\end{fact}

\begin{theorem}[Johnson bound]\label{thm:johnson-interleaved}
For every interleaved word $\orc$ of $\RS[\E,\dom,2^\kappa]$ (rate $\rate$) and every $\slack \in (0, 1-\sqrt{\rate})$: $\ |\lst_{1-\sqrt{\rate}-\slack}(\orc)| \le \frac{1}{2\slack\sqrt{\rate}}$. (see \S\ref{app:johnson} for a self-contained proof)
\end{theorem}

\subsubsection{PCS and round-by-round soundness}\label{sec:iopcs}

\begin{definition}[$L$-list binding, informal]\label{def:listbinding}
A \emph{polynomial commitment scheme} (PCS) is \emph{$L$-list binding} if every commitment transcript $\mathsf{cm}$ determines a set $\polyset_{\mathsf{cm}}$ of at most $L$ multilinear polynomials, not necessarily efficiently computable, such that for every unbounded prover and every collection of opened claims: if no member of $\polyset_{\mathsf{cm}}$ satisfies all the claims, the verifier rejects except with small probability. $L = 1$ is standard binding.
\end{definition}

\begin{definition}[RBR soundness, informal; cf.~\cite{CCHLRR19,ACFY25}]\label{def:rbr}
A public-coin protocol is \emph{round-by-round (RBR) sound with errors} $(\varepsilon_1, \varepsilon_2, \dots)$, one per verifier message, if there is an \emph{invariant} on partial transcripts (not necessarily efficiently computable) marking the positions where a cheating prover has already lost:
\begin{enumerate}[itemsep=1pt,label=(\roman*)]
\item on a false statement, the invariant holds before any message is sent;
\item the prover alone cannot escape it: it survives every prover message;
\item the $i$-th verifier message lets it escape with probability at most $\varepsilon_i$;
\item if it survives to the end of the interaction, the verifier rejects.
\end{enumerate}
\end{definition}

Interactive soundness error is then at most $\sum_i \varepsilon_i$. For a canonical random oracle on the full transcript, the doomed-state sparse-relation argument of~\cite[Definition 5.3 and proof of Theorem 5.8]{CCHLRR18} gives ordinary soundness by a union bound over fresh challenge queries, including verification, with each query charged its round's error. This ROM corollary is not an assertion that every compression call is one such query, nor an application of the different IOP compiler of~\cite{BCS16}. The buffered duplex groups consecutive samples and omits previous challenges from its stored history; \S\ref{sec:e2e-fiat-shamir} gives the necessary ancestor-completion adapter and its explicit query and grouping losses.

%=====================================================================
\subsection{Reed--Solomon codes over binary fields}\label{sec:binary-rs}
%=====================================================================

The encoder $\Enc_{\kappa,\rexp}$ of Definition~\ref{def:enc} evaluates $2^\kappa$ novel-basis coefficients on an additive domain of size $\blen=2^{\kappa+\rexp}$. Here $\rexp\ge1$ and $\kappa+\rexp\le64$. The protocol uses two properties, proved in Annex~\ref{app:ntt}:
\begin{itemize}[itemsep=2pt]
\item \emph{fast encoding, for the prover:} $\Enc_{\kappa,\rexp}$ is computed by the additive NTT with $\tfrac12 \kappa \blen$ multiplications and fewer than $(\kappa + 1)\, \blen$ additions (Proposition~\ref{prop:ntt});
\item \emph{MLE-friendly columns, for the verifier:} for every $x \in \dom$, the weight $W_x : \cube{\kappa} \to \K$, defined by $\Enc_{\kappa,\rexp}(g)[x] = \inner{W_x}{g}$ (the relation being valid for arbitrary $g$), admits an MLE efficiently evaluable anywhere in $\E^{\kappa}$ in $O(\kappa)$ field operations (Lemma~\ref{lem:colweight}).
\end{itemize}

%=====================================================================
\subsection{The protocol}\label{sec:protocol}
%=====================================================================

\begin{protocol}[Whir/Ligerito]\label{fig:protocol}
\ \par\smallskip
\emph{Parameters.} The scheme commits to $f : \cube{\nv} \to \K$ across $\nlev \ge 1$ \emph{levels} $i = 0, \dots, \nlev-1$ plus a plaintext final level $\nlev$: variable counts $\nv = \nv_0 > \nv_1 > \dots > \nv_\nlev \ge 0$ with fold factors $\ell_i := \nv_i - \nv_{i+1} \ge 1$ and $\nv_\nlev$ small (e.g.\ $\nv_\nlev \le 5$); per level, a rate $\rate_i = 2^{-\rexp_i}$, $\rexp_i \ge 1$, giving the encoder $\Enc_i := \Enc_{\kappa_i, \rexp_i}$ of \S\ref{sec:binary-rs} with $\kappa_i := \nv_{i+1}$, i.e.\ the code $\mathcal{C}_i := \RS[\E, \dom_i, 2^{\kappa_i}]$ over its domain $\dom_i \subseteq \K$ of size $\blen_i = 2^{\kappa_i + \rexp_i}$ (requires $\kappa_i + \rexp_i \le 64$); and a query count $\nq_i$. (The proximity radii $\rad_i \in (0, 1 - \sqrt{\rate_i})$ appear only in the analysis.)

\smallskip
\textbf{Commit}$(f : \cube{\nv} \to \K)$: $\mathbf{P}$ sends the oracle $\orc^{(0)}[u, \cdot] := \Enc_0(f(u, \cdot))$, $u \in \cube{\ell_0}$, whose entries are $\K$ elements. No OOD sample is taken.

\smallskip
\textbf{Open}$\bigl(\{(W_t, c_t)\}_{t = 1, \dots, J_0}\bigr)$, $J_0 \ge 1$ weighted claims about $f$: set $f^{(0)} := f$. A round message $h \in \E[X]$, $\deg h \le 2$, answering a claimed sum $c$ travels as $h(0)$ and its $X^2$ coefficient: $\mathbf{V}$ recovers the linear coefficient from $h(0) + h(1) = c$, so that check holds by construction. The sumcheck runs one message ahead: the round-$(j{+}1)$ message is sent right after $\fc_j$, so a level's first message is in hand before its claims are batched. For $i = 0, 1, \dots, \nlev - 1$:
\begin{enumerate}[itemsep=0.5pt,topsep=2pt]
\item \emph{(batch)} $\mathbf{V}$ sends $\lambda \leftarrow \E$; the \emph{batched claim} is $(W, \rc{0}) := \bigl(\sum_{t = 1}^{J_i} \lambda^{t-1} W_t,\; \sum_{t} \lambda^{t-1} c_t\bigr)$. Its first round polynomial $h_1 := \sum_t \lambda^{t-1} h^{(t)}$ is never sent whole: $\mathbf{P}$ sends one \emph{intro polynomial} $h^{(t)}(X) = \sum_{x} \mle{W_t}(X, x)\, \mle{f^{(i)}}(X, x)$ per claim, answering $c_t$, and $\mathbf{V}$ sums them with the powers of $\lambda$. At level $0$, $\mathbf{P}$ sends them after $\lambda$, and with $J_0 = 1$ (the stacked opening, \S\ref{sec:stacking}) no $\lambda$ is drawn. At level $i \ge 1$ they arrive while the claims are made: the residual claim's (which keeps $\lambda^0 = 1$) right after the last fold of level $i-1$, the OOD claim's with $\cood$ in step 3(b), and one for the consistency claims after their openings in step 3(c), these claims being batched among themselves by the next powers of $\lambda$, so that each is one term of the sum.
\item \emph{(sumcheck~\cite{LFKN92})} For $j = 1, \dots, \ell_i$, $\mathbf{P}$ holding $h_j$, honestly $h_j(X) = \sum_{x} \mle{W}(\fc^{<j}, X, x)\, \mle{f^{(i)}}(\fc^{<j}, X, x)$ where $\fc^{<j} := (\fc_1, \dots, \fc_{j-1})$, with $h_j(0) + h_j(1) = \rc{j-1}$: $\mathbf{V}$ samples $\fc_j \leftarrow \E$ and sends it to the prover; both set $\rc{j} := h_j(\fc_j)$; $\mathbf{P}$ sends $h_{j+1}$, the round-$(j{+}1)$ polynomial of the claim $(\mle{W}(\fc^{\le j}, \cdot), \rc{j})$, which for $j = \ell_i$ is the residual claim's intro polynomial at the next level.\\
Set $\fc := (\fc_1, \dots, \fc_{\ell_i})$, $c' := \rc{\ell_i}$, and, for $x \in \cube{\nv_{i+1}}$: \ $f^{(i+1)}(x) := \mle{f^{(i)}}(\fc, x)$, \ $W'(x) := \mle{W}(\fc, x)$. The \emph{residual claim} $(W', c')$ is a weighted claim about the folded multilinear $f^{(i+1)} : \cube{\nv_{i+1}} \to \E$.
\item If $i < \nlev - 1$:
\begin{enumerate}[itemsep=0.5pt,topsep=1pt,label=(\alph*)]
\item \emph{(commit)} $\mathbf{P}$ sends the oracle $\orc^{(i+1)}[u, \cdot] := \Enc_{i+1}(f^{(i+1)}(u, \cdot))$, $u \in \cube{\ell_{i+1}}$.
\item \emph{(ood)} $\mathbf{V}$ sends $z \leftarrow \E^{\nv_{i+1}}$; $\mathbf{P}$ sends $\cood \in \E$ (honestly $\mle{f^{(i+1)}}(z)$), then the intro polynomial of the claim $(\eq(z, \cdot), \cood)$.
\item \emph{(query)} $\mathbf{V}$ sends $\nq_i$ stratified columns $\qset_i = (x_q)_{q = 1, \dots, \nq_i}$ of $\dom_i$ (Definition~\ref{def:strata}), then the next level's $\lambda$; it reads $\orc^{(i)}[\cdot, x_q]$ and computes $c_q := \sum_{u} \eq(\fc, u)\, \orc^{(i)}[u, x_q] = \orc^{(i)}_\fc[x_q]$; honestly $c_q = \Enc_i(f^{(i+1)})[x_q]$, by Fact~\ref{fact:fold-decode}. $\mathbf{P}$ sends the intro polynomial of the consistency claims' batch.
\item Claims for level $i+1$ (so $J_{i+1} := \nq_i + 2$), in batching order: the residual claim $(W', c')$, the OOD claim $(\eq(z, \cdot), \cood)$, and the $\nq_i$ \emph{consistency claims} $(W^{(i)}_{x_q}, c_q)$, where $W^{(i)}_x$ are the column weights of $\Enc_i$ (\S\ref{sec:binary-rs}): the $q$-th asserts $\Enc_i(f^{(i+1)})[x_q] = c_q$.
\end{enumerate}
\item If $i = \nlev - 1$: \emph{(final)} $\mathbf{P}$ sends the multilinear $\ffinal : \cube{\nv_\nlev} \to \E$; $\mathbf{V}$ sends $\nq_{\nlev-1}$ stratified columns $\qset_{\nlev-1}$ of $\dom_{\nlev-1}$ and computes $c_q$ from $\orc^{(\nlev-1)}$ as in 3(c). Consistency is again a weighted claim, $\Enc_{\nlev-1}(\ffinal)[x_q] = \inner{W^{(\nlev-1)}_{x_q}}{\ffinal}$ (Lemma~\ref{lem:colweight}), so the level closes like the others: $\mathbf{V}$ samples and sends $\lambda \leftarrow \E$, batching the $\nq_{\nlev-1}$ consistency claims with the residual claim $(W', c')$ into one claim about $\ffinal$, whose first round polynomial is the residual claim's intro polynomial plus $\lambda$ times the one $\mathbf{P}$ sends for the consistency claims; $\nv_\nlev$ further sumcheck rounds discharge it, and $\mathbf{V}$ closes by evaluating $\mle{\ffinal}$ at their point from the table it holds.
\end{enumerate}
$\mathbf{V}$ accepts iff every check above passed.
\end{protocol}

\begin{definition}[Stratified columns]\label{def:strata}
Number the columns of $\dom_i$ from $0$ to $\blen_i - 1$ in the order of the Merkle tree's leaves, $\blen_i$ a power of two, so that a column's top bits name the subtree holding its leaf. The $\nq_i$ columns of a query message are drawn in groups, one for each binary digit one of $\nq_i$, the largest first. A group of $2^{\strh}$ columns fixes the top $\strs := \min(\strh, \log_2 \blen_i)$ bits of its $j$-th column, $0 \le j < 2^{\strh}$, to $j \bmod 2^{\strs}$, and draws its other bits uniformly, every column independently. Each of the $2^{\strs}$ cosets of a group's fixed bits thus holds $2^{\strh - \strs}$ of its columns, each uniform in the coset.
\end{definition}

\begin{lemma}\label{lem:strata}
If a set of columns has fewer than $(1 - \rad_i)\blen_i$ elements, all $\nq_i$ stratified columns of a query message lie in it with probability less than $(1 - \rad_i)^{\nq_i}$, as i.i.d.\ uniform columns would.
\end{lemma}

\begin{proof}
Take one group, of $2^{\strh}$ columns, and let $\stcov{1}, \dots, \stcov{2^{\strs}}$ be the fractions of the cosets of its fixed bits that the set covers, whose mean is the set's density, below $1 - \rad_i$. The group lies in the set with probability $(\stcov{1} \cdots \stcov{2^{\strs}})^{2^{\strh - \strs}} \le \bigl(2^{-\strs}(\stcov{1} + \dots + \stcov{2^{\strs}})\bigr)^{2^{\strh}} < (1 - \rad_i)^{2^{\strh}}$, by the inequality of arithmetic and geometric means. The groups are independent and their sizes sum to $\nq_i$.
\end{proof}

The verifier's Merkle check gains from the strata: the largest group's columns reach every node of the tree's top $\strs$ levels, a complete subtree, which is hashed once, and every column's own path stops where its fixed bits begin (Annex~\ref{annex:rec}, \S\ref{rec:verifier}).

All weights in the protocol are MLE-friendly and multilinear, so every sumcheck summand $\mle{W} \cdot \mle{f^{(i)}}$ has individual degree at most $2$. At level $0$ that summand pairs an $\E$-valued weight with the $\K$-valued $f$; every later level is $\E$ throughout.

The analysis below reads a level's first round polynomial $h_1$ as one message: every part of it is fixed before $\fc_1$ is drawn, and each part answers its own claim, so $h_1(0) + h_1(1) = \rc{0}$.

Completeness is perfect: the honest prover passes every check ($h_j(0) + h_j(1) = \rc{j-1}$ by definition of $h_j$, $c_q = \Enc_i(f^{(i+1)})[x_q]$ by Fact~\ref{fact:fold-decode}, the OOD and residual claims by construction, the final checks as identities).

%=====================================================================
\subsection{Round-by-round soundness}\label{sec:soundness}
%=====================================================================

Fix a commitment $\mathsf{cm} = \orc^{(0)}$ and input claims $\{(W_t, c_t)\}_t$, and define the relation
\[
\relopen \;:=\; \Bigl\{ \bigl(\orc^{(0)}, \{(W_t, c_t)\}_t\bigr) \;:\; \exists\, U \in \lst_{\rad_0}(\orc^{(0)}) \text{ with } \inner{W_t}{g_U} = c_t \text{ for all } t \Bigr\}.
\]
The commitment binds the prover to the list $\polyset_{\mathsf{cm}} := \{\mle{g_U} : U \in \lst_{\rad_0}(\orc^{(0)})\}$, of size at most $L_0 := \tfrac{1}{2\slack_0\sqrt{\rate_0}}$ for $\rad_0 = 1 - \sqrt{\rate_0} - \slack_0$, by the Johnson bound (Theorem~\ref{thm:johnson-interleaved}). An outer protocol that draws challenges between the commitment and the opening must therefore have each of them hold against every member of this list, at a cost of a factor $L_0$ on its error (\S\ref{sec:e2e-ledger}).

\begin{theorem}[RBR soundness]\label{thm:rbr}
Suppose that for every level $i$: $\rad_i = 1 - \sqrt{\rate_i} - \slack_i$ with $\slack_i \in (0, 1-\sqrt{\rate_i})$, $L_i := \tfrac{1}{2\slack_i\sqrt{\rate_i}}$, and $\epsilon_i := \mcanumi{i}/|\E|$ is the mutual correlated agreement error of Theorem~\ref{thm:mca-johnson} below for $\mathcal{C}_i$ at radius $\rad_i$. Then the Protocol~\ref{fig:protocol} is a $L_0$-list binding PCS, and its opening phase $\relopen$ has the following RBR soundness error at each verifier message:
\[
\begin{array}{llcl}
\textnormal{batching, level } i: & \varepsilon^{\mathrm{batch}}_i &=& \dfrac{(J_i - 1)\, L_i}{|\E|} , \\[2ex]
\textnormal{fold challenge } \fc_j, \textnormal{ level } i: & \varepsilon^{\mathrm{fold}}_{i,j} &=& \dfrac{2 L_i}{|\E|} + \epsilon_i , \\[2ex]
\textnormal{OOD challenge entering level } i \ge 1: & \varepsilon^{\mathrm{ood}}_i &=& \dbinom{L_i}{2} \cdot \dfrac{\nv_i}{|\E|} , \\[2ex]
\textnormal{query message } \qset_i: & \varepsilon^{\mathrm{query}}_i &=& (1 - \rad_i)^{\nq_i} , \\[2ex]
\textnormal{final batching challenge } \lambda: & \varepsilon^{\mathrm{fin}} &=& \dfrac{\nq_{\nlev-1}}{|\E|} , \\[2ex]
\textnormal{each of the } \nv_\nlev \textnormal{ closing rounds}: & \varepsilon^{\mathrm{tail}} &=& \dfrac{2}{|\E|} .
\end{array}
\]
Here $J_0$ is the number of input claims and $J_i = \nq_{i-1} + 2$ for $i \ge 1$. The \emph{RBR error}, which governs the Fiat--Shamir compilation (\S\ref{sec:iopcs}), is the maximum of these entries over all rounds.
\end{theorem}

\paragraph{The ring switch's place.} In leanVM the input claims come from ring switching (Annex~\ref{annex:rs}), whose challenges $\rsgamma$ and $\Phi$ are drawn before level $0$'s batching challenge. Its errors are fixed by the committed polynomial, which the commitment binds only to a list, so Lemma~\ref{lem:rs-family} is unioned over that list: $(\nrs - 1) L_0/|\E|$ for $\rsgamma$ and $2^{32} L_0/|\E|$ for $\Phi$, whose discrepancy has total degree below $2^{32}$ (\S\ref{rs:map}). These are level-$0$ terms, entered with $\varepsilon^{\mathrm{batch}}_0$. A deeper level's oracle is sent after every ring-switch challenge, and its claims are those of step 3(d), so its batching term is $(J_i - 1) L_i/|\E|$ alone.

The remainder of this section proves the theorem: \S\ref{sec:tools} states the coding-theoretic tools, \S\ref{pcs:proof} the invariant and the round-by-round argument.

\subsubsection{Proximity tools}\label{sec:tools}

The engine is \emph{mutual correlated agreement} (MCA)~\cite{ACFY25}. For a two-row word $\orc$, recall $\orc_\fc = (1-\fc)\orc[0,\cdot] + \fc\,\orc[1,\cdot]$.\footnote{\cite{BCHKS25} state the results for the affine line $\orc[0,\cdot] + \fc\,\orc[1,\cdot]$; in characteristic $2$, $\orc_\fc = \orc[0,\cdot] + \fc(\orc[0,\cdot] + \orc[1,\cdot])$, and this invertible change of pair maps codeword pairs to codeword pairs, so the formulations are equivalent.}

\begin{theorem}[MCA in the unique-decoding regime; adapted from {\cite[Cor.~1.4]{BCHKS25}}]\label{thm:mca-udr}
Let $\RS[\E,\dom,\kdim]$ have block length $\blen = |\dom|$ and relative minimum distance $\dmin := 1 - (\kdim-1)/\blen$. Suppose
\[
\dmin \ge 3\sqrt{\frac{2}{\blen}}
\qquad\text{and}\qquad
\rad \in \left[\frac{\dmin}{3},\,\frac{\dmin}{2} - \frac{3}{\dmin \blen}\right].
\]
Then, for every two-row word $\orc$,
\[
\Pr_{\fc \leftarrow \E}\!\left[\; \exists \aset \subseteq \dom,\ |\aset| \ge (1-\rad) \blen :\ \begin{array}{l} \orc_\fc \text{ agrees with the code on } \aset, \text{ and} \\ \orc \text{ does not agree with the interleaved code on } \aset \end{array} \right]
\;\le\; \frac{\rad \blen + 1}{|\E|}.
\]
\end{theorem}

We record Theorem~\ref{thm:mca-udr} only for comparison, we don't need it in the List-Decoding Regime, we instead need:

\begin{theorem}[MCA up to the Johnson bound; {\cite[Thm.~4.6]{BCHKS25}}]\label{thm:mca-johnson}
Let $\RS[\E,\dom,\kdim]$ have block length $\blen = |\dom|$ and set $\rate := (\kdim-1)/\blen$, the \emph{slightly reduced rate} (\cite{BCHKS25} parametrize by the degree bound $\kdim - 1$, not the dimension $\kdim$). Let $\rad \in (0, 1-\sqrt{\rate})$, and set $\slack := 1 - \sqrt{\rate} - \rad$ and $\mcapar := \max\{\lceil \tfrac{\sqrt{\rate}}{\slack} \rceil, 3\}$ (denoted $m$ in~\cite{BCHKS25}). For every two-row word $\orc$,
\[
\Pr_{\fc \leftarrow \E}\!\left[\; \exists \aset \subseteq \dom,\ |\aset| \ge (1-\rad) \blen :\ \begin{array}{l} \orc_\fc \text{ agrees with the code on } \aset, \text{ and} \\ \orc \text{ does not agree with the interleaved code on } \aset \end{array} \right] \;\le\; \frac{\mcanum}{|\E|},
\]
where
\[
\mcanum \;=\; \frac{2(\mcapar + \tfrac12)^5 + 3(\mcapar + \tfrac12)\rad\rate}{3\rate^{3/2}} \cdot \blen \;+\; \frac{\mcapar + \tfrac12}{\sqrt{\rate}}
\]
\end{theorem}

We write $\epsilon := \mcanum/|\E|$ for this MCA error at radius $\rad$. MCA is exactly what Lemma~\ref{lem:fold-list} needs: folding should not create \emph{new} nearby codewords.

\begin{lemma}[Folding preserves lists]\label{lem:fold-list}
Let $\RS[\E,\dom,\kdim]$ have rate $\rate$, let $\rad \in (0, 1-\rate)$, and let $\epsilon$ be an MCA error at radius $\rad$, i.e.\ a bound on the probability of Theorem~\ref{thm:mca-johnson} for every two-row word. For every interleaved word $\orc \in \E^{\cube{\ell} \times \dom}$, $\ell \ge 1$:
\[
\Pr_{\fc \leftarrow \E}\Bigl[\, \lst_\rad(\orc_\fc) \neq \{ U_\fc : U \in \lst_\rad(\orc) \} \,\Bigr] \;\le\; \epsilon .
\]
\end{lemma}

The bound does not grow with $\ell$. This is a theorem of Jo~\cite[Thm.~4.4]{Jo26}: for a code linear over a field $\F$, and a fold whose two coefficients are functions of a seed drawn uniformly from a set of at most $|\F|$ elements, the MCA error of every row interleaving of the code equals that of the code itself, at every radius. Both hypotheses hold here, with $\F = \E$:
\begin{itemize}[itemsep=1pt]
\item \emph{The code is $\E$-linear.} Its domain lies in $\K$, but $\RS[\E,\dom,\kdim]$ is an $\E$-subspace of $\E^\dom$, and so is each of its row interleavings.
\item \emph{The seed set is no larger than the field.} The seed is the fold challenge $\fc$, drawn from $\E$ itself, and the coefficients are $(1-\fc, \fc)$.
\end{itemize}
The theorem holds at every radius, hence in the list-decoding regime of Theorem~\ref{thm:mca-johnson}, and the event it bounds is the MCA event of that theorem, with one agreement set shared by the fold and by every row. A row-by-row union bound would instead pay $2^{\ell-1}\epsilon$.

\begin{proof}
$\supseteq$ holds always: $U \in \lst_\rad(\orc)$ agreeing with $\orc$ on $\aset$ gives $U_\fc$ agreeing with $\orc_\fc$ on $\aset$. For $\subseteq$: read $\orc$ as the two-row word $(\orc[(0,\cdot),\cdot], \orc[(1,\cdot),\cdot])$ over the $2^{\ell-1}$-fold row interleaving of the code, whose fold at $\fc$ is $\orc_\fc$, and consider its MCA event with respect to that interleaved code. By \cite[Thm.~4.4]{Jo26}, that event has probability at most the MCA error of a two-row word of $\RS[\E,\dom,\kdim]$, hence at most $\epsilon$. Suppose it does not occur, and let $V \in \lst_\rad(\orc_\fc)$ with common agreement set $\aset$, $|\aset| \ge (1-\rad)\blen$. Then $\orc_\fc$ agrees on $\aset$ with the interleaved codeword $V$, so (as the event failed) there is an interleaved codeword $W$ whose rows $W[(\iota,u),\cdot]$ equal $\orc[(\iota,u),\cdot]$ on $\aset$ for every $\iota \in \{0,1\}$ and $u \in \cube{\ell-1}$; $W$ lies in $\lst_\rad(\orc)$. Finally $W_\fc[u,\cdot]$ agrees with $V[u,\cdot]$ on $\aset$, and $|\aset| \ge (1-\rad)\blen > \kdim$ while distinct codewords agree on fewer than $\kdim$ points, so $V = W_\fc$.
\end{proof}

\begin{lemma}[OOD separation]\label{lem:ood}
Let $\orc$ be an interleaved word with $|\lst_\rad(\orc)| \le L$, decoding to $\nv$-variable multilinears. Then $\Pr_{z \leftarrow \E^{\nv}}[\exists \text{ distinct } U, U' \in \lst_\rad(\orc) : \mle{g_U}(z) = \mle{g_{U'}}(z)] \le \binom{L}{2} \nv / |\E|$.
\end{lemma}

\begin{proof}
Distinct interleaved codewords have distinct multilinear extensions; apply Lemma~\ref{lem:sz} to each difference and union bound.
\end{proof}

\subsubsection{Proof of the main theorem}\label{pcs:proof}

At level $i$, abbreviate $\fc^{\le j} := (\fc_1, \dots, \fc_j)$ and $\fc^{<j} := \fc^{\le j-1}$; thus $\fc^{\le 0}$ is empty and $\orc^{(i)}_{\fc^{\le 0}} = \orc^{(i)}$. For $0 \le j \le \ell_i$, the \emph{round-$j$ claim} is the weighted claim
\[
\bigl(\, \mle{W}(\fc^{\le j}, \cdot),\; \rc{j} \,\bigr) \qquad \text{about multilinears } g : \cube{\nv_i - j} \to \E,
\]
where $W$, $\rc{0}$, and $\rc{j} = h_j(\fc_j)$ are read off the transcript as in Protocol~\ref{fig:protocol}; an interleaved codeword satisfies a claim if its decoded multilinear does. The round-$0$ claim is the batched claim, and the round-$\ell_i$ claim is the residual claim $(W', c')$.

We now define the invariant required by Definition~\ref{def:rbr}, a condition on the partial transcripts of the opening. In one phrase: \emph{every codeword close to the current fold violates the current claim}. Precisely, the list below gives the condition just after each verifier message; a transcript ending with prover messages inherits the condition of its last verifier message. Every condition also includes one implicit disjunct, the \emph{escape clause}: \emph{``some verifier check has already failed''} (such transcripts are rejected no matter what follows).
\begin{itemize}[itemsep=1pt]
\item \emph{Start of level $i$} (before any message for $i = 0$, after the query message $J_{i-1}$ otherwise): every $U \in \lst_{\rad_i}(\orc^{(i)})$ violates at least one of the level's $J_i$ claims (the input claims for $i = 0$, the claims of step 3(d) of Protocol~\ref{fig:protocol} otherwise).
\item \emph{After the batching challenge $\lambda$ of level $i$, and after each of its fold challenges $\fc_j$:} every $U \in \lst_{\rad_i}(\orc^{(i)}_{\fc^{\le j}})$ violates the round-$j$ claim ($\lambda$ being the case $j = 0$).
\item \emph{After the OOD challenge $z$ closing level $i$ ($i < \nlev-1$):} the round-$\ell_i$ instance above holds, and the multilinear extensions decoded from $\lst_{\rad_{i+1}}(\orc^{(i+1)})$ are pairwise distinct at $z$.
\item \emph{After the final query message $\qset_{\nlev-1}$:} $\ffinal$ violates one of the level's $\nq_{\nlev-1}+1$ claims.
\item \emph{After the final $\lambda$, and after each of the $\nv_\nlev$ rounds that follow (complete transcript):} $\ffinal$ violates the running claim. The terminal one $\mathbf{V}$ checks against $\mle{\ffinal}$ itself, so it rejects.
\end{itemize}
Conditions (i), (ii), (iv) of Definition~\ref{def:rbr} hold by inspection: on a false statement, the first item is exactly the failure of $\relopen$; a prover message changes nothing an instance depends on, except possibly failing a check, which only switches the escape clause on; and the last item is precisely ``some check failed'', so such a complete transcript is rejected. The separation clause of the OOD instance is there because the query message consumes it: it pins the new oracle's list to at most one candidate consistent with the prover's answer. It remains to bound the escape probabilities (iii); in the proof below we assume all checks so far pass, since otherwise the escape clause already holds.

\begin{proof}[Proof of Theorem~\ref{thm:rbr}]
If the statement lies outside $\relopen$, then before any message is sent every $U \in \lst_{\rad_0}(\orc^{(0)})$ violates at least one input claim, so the invariant holds. We follow the protocol and bound, for each verifier message, the probability that the invariant fails to carry over. Every list below consists of interleaved words of $\mathcal{C}_i$ within radius $\rad_i$, hence has size at most $L_i$ by Theorem~\ref{thm:johnson-interleaved}.

\medskip\noindent\emph{Batching challenge $\lambda$.}
Assume that every $U \in \lst_{\rad_i}(\orc^{(i)})$ violates at least one of the $J_i$ claims $(W_t, c_t)$. Fix such a $U$. Its discrepancy against the batched claim is $\sum_{t} \lambda^{t - 1} d_t$ with $d_t := \inner{W_t}{g_U} - c_t$, a polynomial in $\lambda$ of degree less than $J_i$, nonzero because the $d_t$ are not all zero. By Lemma~\ref{lem:sz} it vanishes with probability at most $(J_i - 1)/|\E|$, so a union bound over the at most $L_i$ list members leaves every $U$ violating the round-$0$ claim except with probability $(J_i - 1) L_i / |\E|$.

\medskip\noindent\emph{Fold challenge $\fc_j$.}
Assume that every $U \in \lst := \lst_{\rad_i}(\orc^{(i)}_{\fc^{<j}})$ violates the round-$(j{-}1)$ claim. Since all checks so far pass (otherwise the verifier trivially rejects), $h_j(0) + h_j(1) = \rc{j-1}$. For $U \in \lst$, let
\[
H_U(X) \;:=\; \sum_{x \in \cube{\nv_i - j}} \mle{W}(\fc^{<j}, X, x)\, \mle{g_U}(X, x)
\]
be the honest round polynomial for $U$; it has degree at most $2$, since $\mle{W}$ and $\mle{g_U}$ are multilinear. Summing over $X \in \{0, 1\}$ turns the multilinear evaluations back into values, so
\[
H_U(0) + H_U(1) \;=\; \bigl\langle\, \mle{W}(\fc^{<j}, \cdot),\; g_U \,\bigr\rangle \;\neq\; \rc{j-1} \;=\; h_j(0) + h_j(1),
\]
the inequality being the assumed violation; hence $h_j \neq H_U$ as polynomials. Two bad events over $\fc_j$:
\begin{itemize}[itemsep=1pt]
\item $B_1$: $h_j(\fc_j) = H_U(\fc_j)$ for some $U \in \lst$. For each $U$, two distinct polynomials of degree at most $2$ agree on at most $2$ points, so $\Pr[B_1] \le 2 L_i / |\E|$.
\item $B_2$: $\lst_{\rad_i}(\orc^{(i)}_{\fc^{\le j}}) \neq \{ U_{\fc_j} : U \in \lst \}$, i.e.\ folding does not map the old list onto the new one. Lemma~\ref{lem:fold-list}, applied to the $2^{\ell_i - j + 1}$-row matrix $\orc^{(i)}_{\fc^{<j}}$, gives $\Pr[B_2] \le \epsilon_i$.
\end{itemize}
If neither occurs, every $V \in \lst_{\rad_i}(\orc^{(i)}_{\fc^{\le j}})$ is $V = U_{\fc_j}$ for some $U \in \lst$, with $\mle{g_V} = \mle{g_U}(\fc_j, \cdot)$ by Fact~\ref{fact:fold-decode}, so
\[
\bigl\langle\, \mle{W}(\fc^{\le j}, \cdot),\; g_V \,\bigr\rangle
\;=\; \sum_{x \in \cube{\nv_i - j}} \mle{W}(\fc^{<j}, \fc_j, x)\, \mle{g_U}(\fc_j, x)
\;=\; H_U(\fc_j) \;\neq\; h_j(\fc_j) \;=\; \rc{j} ,
\]
using $\neg B_1$ for the inequality: every $V$ violates the round-$j$ claim. The total error is $2 L_i/|\E| + \epsilon_i$.

\medskip\noindent\emph{The level interface.}
After the last fold challenge of level $i < \nlev - 1$, we know: every codeword within $\rad_i$ of the single row $c := \orc^{(i)}_\fc \in \E^{\dom_i}$ violates the residual claim $(W', c')$. The prover now sends the new oracle $\orc^{(i+1)}$.

\medskip\noindent\emph{OOD challenge $z$.}
After $z$, two properties must hold: the one above, which does not involve $z$; and pairwise distinctness at $z$ of the multilinear extensions decoded from $\lst_{\rad_{i+1}}(\orc^{(i+1)})$, a list already determined by the prover's message. The latter fails with probability at most $\binom{L_{i+1}}{2} \nv_{i+1} / |\E|$ over $z$, by Lemma~\ref{lem:ood}.

\medskip\noindent\emph{Query message $\qset_i$.}
Assume both properties of the previous step; the prover has since sent its answer $\cood$. The values $\mle{g_U}(z)$ are pairwise distinct across $\lst_{\rad_{i+1}}(\orc^{(i+1)})$, so at most one member $U^\ast$ can satisfy $\mle{g_{U^\ast}}(z) = \cood$: every other member violates the OOD claim $(\eq(z, \cdot),\, \cood)$, whatever $\qset_i$ is. If no such $U^\ast$ exists, we are done. Otherwise write $g^\ast := g_{U^\ast}$, a multilinear on $\cube{\nv_{i+1}}$, and distinguish two cases.
\begin{itemize}[itemsep=1pt]
\item \emph{$\Enc_i(g^\ast)$ agrees with $c$ on at least $(1 - \rad_i) \blen_i$ coordinates.} Then the codeword $\Enc_i(g^\ast)$ lies in $\lst_{\rad_i}(c)$, hence violates the residual claim by assumption: $\inner{W'}{g^\ast} \neq c'$, whatever $\qset_i$ is.
\item \emph{Otherwise.} $\Enc_i(g^\ast)$ and $c$ agree on fewer than $(1 - \rad_i)\blen_i$ columns, a set fixed before $\qset_i$ is drawn, so all $\nq_i$ stratified queries lie in it with probability less than $(1 - \rad_i)^{\nq_i}$ (Lemma~\ref{lem:strata}); and any disagreeing query $x_q$ makes $U^\ast$ violate the consistency claim $(W^{(i)}_{x_q},\, c_q)$, since $c_q = c[x_q]$.
\end{itemize}
Except with probability $(1 - \rad_i)^{\nq_i}$, every member of $\lst_{\rad_{i+1}}(\orc^{(i+1)})$ thus violates at least one of the $J_{i+1}$ claims, which is exactly the assumption of the batching step at level $i + 1$.

\medskip\noindent\emph{Final level.}
At level $\nlev - 1$ the fold challenges end as above: every codeword within $\rad_{\nlev-1}$ of $c := \orc^{(\nlev-1)}_\fc$ violates the residual claim. The prover sends $\ffinal$. If $\Enc_{\nlev-1}(\ffinal)$ agrees with $c$ on at least $(1 - \rad_{\nlev-1}) \blen_{\nlev-1}$ coordinates then $\Enc_{\nlev-1}(\ffinal) \in \lst_{\rad_{\nlev-1}}(c)$, so $\ffinal$ violates the residual claim; otherwise, except with probability $(1 - \rad_{\nlev-1})^{\nq_{\nlev-1}}$ (Lemma~\ref{lem:strata}), some query disagrees and $\ffinal$ violates that consistency claim. Either way it violates one of the $\nq_{\nlev-1}+1$ batched claims, hence the batch itself except with probability $\nq_{\nlev-1}/|\E|$ (Lemma~\ref{lem:sz}), with no union over a list: $\ffinal$ is transmitted, not one of the codewords near an oracle. The $\nv_\nlev$ rounds that follow are a plain sumcheck (Fact~\ref{fact:sumcheck}) on a polynomial that $\mathbf{V}$ holds, so the terminal check against $\mle{\ffinal}$ fails.

\end{proof}


%=====================================================================
\subsection{Optimizations}\label{sec:pcs-opt}
%=====================================================================

The committed $f$ is the stack of \S\ref{sec:stacking}: $\nstack$ field elements in $\K$, then zeros up to $2^{\nv}$. Let the row index be the stack's most significant $\ell_0$ coordinates: $f(u, x) := q(x, u)$ for $u \in \cube{\ell_0}$ and $x \in \cube{\nv - \ell_0}$, \S\ref{app:ntt}. Row $u$ is then the block of $2^{\nv - \ell_0}$ consecutive elements of $q$ at offset $\langle u \rangle 2^{\nv - \ell_0}$, the padding is whole rows, and only $\nrows=\lceil\nstack/2^{\nv-\ell_0}\rceil$ of the $2^{\ell_0}$ rows hold data. So we get the following optimizations at commitment:
\begin{itemize}[itemsep=1pt,topsep=2pt]
\item the FFT-encoder only runs on $\nrows$ rows (and not the full $2^{\ell_0}$), since FFT-ing zeros gives zeros;
\item a Merkle leaf only contains $\nrows$ elements of $\K$ (proof size optimization), we keep the convention of hashing the full $2^{\ell_0}$ values for simplicity, but we hash them in reverse order, starting from the padding zeros;
\item the prover can thus preprocess the internal state of BLAKE2s after hashing those zeros, and then hash the useful part of each leaf starting from this state.
\end{itemize}

Conclusion: at commitment, both the FFT and the Merkle leaf hashing cost are proportional to $\nstack$ (not its next power of two). Unfortunately we still pay for the padding inside the $\nstack$ field elements, because some instruction counts are not powers of two. This could be prevented using the Jagged PCS \cite{jagged_pcs}.

%=====================================================================
\subsection{The hiding commitment}\label{sec:pcs-zk}
%=====================================================================

A leaf proved in zero knowledge commits to its stack and opens it with the hiding variant of Protocol~\ref{fig:protocol} below, after VEIL~\cite{DHRR26}: every committed row of $\orc^{(0)}$ is padded past its power of two, and one committed row is uniformly random. Only level $0$ changes; every deeper level commits an unpadded fold and runs as before. Write $\kappa := \kappa_0$ and $\rexp := \rexp_0$, so that $\dom_0 = \subsp_{\kappa + \rexp}$, and set
\[
\npad := \nq_0 + 2, \qquad \padshift := e_{\kappa + \rexp},
\]
the first basis element past $\dom_0$.

\paragraph{Commit.} For each committed row $u$ the prover draws a uniform padding $\padc_u \in \K^{\npad}$ and sends
\[
\orc^{(0)}[u, \cdot] := \Bigl(\mpoly_{f(u, \cdot)}(x) + \Wn{\kappa}(x + \padshift) \sum_{j < \npad} \padc_u[j]\, \novel_j(x)\Bigr)_{x \in \dom_0}.
\]
Its last committed row $\rlane$, $\langle \rlane \rangle = \nrows - 1$, holds uniform elements of $\K$ and is padded like the others; no input claim's weight touches it.
\begin{itemize}[itemsep=1pt,topsep=2pt]
\item \emph{The code.} $\Wn{\kappa}$ is $\Ftwo$-linear (Lemma~\ref{lem:subspace}), so $\Wn{\kappa}(X + \padshift) = \Wn{\kappa}(X) + \Wn{\kappa}(\padshift)$, and $\Wn{\kappa} \novel_j = \novel_{2^\kappa + j}$ for $j < 2^\kappa$. A row is therefore the encoding of $2^\kappa + \npad$ novel-basis coefficients ($\npad \le 2^\kappa$), a codeword of $\cpad_0 := \RS[\E, \dom_0, 2^\kappa + \npad]$ on the same domain. The additive NTT computes it with one more butterfly layer, the one at layer $\rexp - 1$ that pairs the message with the padding: the replica of the message on the coset $c + \subsp_\kappa$ of $\dom_0$ starts as the message plus $\Wn{\kappa}(c + \padshift)$ times the padding, which the encoder adds as its first pass reads the message.
\item \emph{No point is left unpadded.} $\Wn{\kappa}$ vanishes on $\subsp_\kappa$, the first $2^\kappa$ points of $\dom_0$, so a padding by $\novel_{2^\kappa + j} = \Wn{\kappa} \novel_j$ alone would leave every query landing there unmasked. $\Wn{\kappa}$ maps $\dom_0$ onto the span of $\Wn{\kappa}(e_\kappa), \dots, \Wn{\kappa}(e_{\kappa + \rexp - 1})$ and is injective on $\langle e_\kappa, \dots, e_{63} \rangle$, so $\Wn{\kappa}(\padshift)$ lies outside that span, and $\Wn{\kappa}(x + \padshift) \neq 0$ for every $x \in \dom_0$.
\item \emph{Decoding.} A polynomial of degree below $2^\kappa + \npad$ is uniquely $\mpoly_g + \Wn{\kappa}(X + \padshift) R$ with $\deg R < \npad$ (division by $\Wn{\kappa}(X + \padshift)$, of degree $2^\kappa$), and a row $U$ of a codeword of $\cpad_0$ decodes to that $g$. The claims are on $g$; the padding is no part of the committed polynomial.
\end{itemize}

\paragraph{Open.} Right after the last lane-fold challenge $\fc_{\ell_0}$, with $c' = \rc{\ell_0}$ fixed and before the next round polynomial:
\begin{enumerate}[itemsep=1pt,topsep=2pt]
\item In the zero-knowledge proof every scalar the prover sent so far travelled under a one-time key, and $\mathbf{V}$ holds $c'$ only as an expression of keyed values. $\mathbf{P}$ sends $c'$ in the clear, $\mathbf{V}$ checks it against that expression and continues from it, and every later scalar travels in the clear.
\item $\mathbf{P}$ sends the paddings folded as the rows are, $\padfold := \sum_{u} \eq(\fc, u)\, \padc_u \in \E^{\npad}$.
\item At level $0$, step 3(c) takes the padding off each folded column:
\[
c_q := \sum_{u} \eq(\fc, u)\, \orc^{(0)}[u, x_q] + \Wn{\kappa}(x_q + \padshift) \sum_{j < \npad} \padfold[j]\, \novel_j(x_q)
\]
(the subtraction, in characteristic $2$), which is $\Enc_0(f^{(1)})[x_q]$ for an honest prover. $\mathbf{V}$ evaluates $\Wn{\kappa}(x_q)$ and the $\novel_j(x_q)$ by the recurrence of Lemma~\ref{lem:subspace}, adds the constant $\Wn{\kappa}(\padshift)$, and folds the sum one bit of $j$ at a time: $\kappa$ squarings and fewer than $\npad$ products per query.
\end{enumerate}
Everything else is Protocol~\ref{fig:protocol}; the induced level-$1$ claims are on the $2^\kappa$ coefficients of $f^{(1)}$ as before.

\paragraph{Soundness.} Theorem~\ref{thm:rbr} holds with level $0$'s code $\cpad_0$, of dimension $2^\kappa + \npad$ and rate $\rate_0 = (2^\kappa + \npad)\, 2^{-\kappa - \rexp}$, in every level-$0$ term: $L_0$, $\epsilon_0$, $\rad_0$ and the query term. The deeper levels keep theirs. The proof reads Protocol~\ref{fig:protocol} with $\cpad_0$ at level $0$ and changes in one place, the query step of level $0$, where $c$ becomes the corrected word $c^\star := \orc^{(0)}_\fc + \Wn{\kappa}(\cdot + \padshift)\, \mpoly_{\padfold}$, $\mpoly_{\padfold} := \sum_{j < \npad} \padfold[j] \novel_j$. If $\Enc_0(g^\ast)$ agrees with $c^\star$ on at least $(1 - \rad_0)\blen_0$ columns, then $\Enc_0(g^\ast) + \Wn{\kappa}(\cdot + \padshift)\, \mpoly_{\padfold}$ is a codeword of $\cpad_0$ within $\rad_0$ of $\orc^{(0)}_\fc$ that decodes to $g^\ast$, which therefore violates the residual claim; otherwise all $\nq_0$ queries land in the agreement set with probability less than $(1 - \rad_0)^{\nq_0}$, as before. $\padfold$ and $c'$ are prover messages and add no error term, and the random row is one more row of the stack, with no claim on it.

The query counts are rederived by the analysis of Theorem~\ref{thm:rbr} for every size and rate, with level $0$'s rate the padded one, at the fixed point $\npad = \nq_0 + 2$: the derivation starts from the unpadded $\nq_0 + 2$ and rederives until $\nq_0$ no longer moves, which it does since a larger padding never needs fewer queries. They differ from the unpadded counts only at rate $1/2$, by a few level-$0$ queries, and most at the smallest stack. The padded code's list bound $L_0$, rederived with its own slack, stays within the $2^{12}$ every challenge before the opening pays (\S\ref{sec:e2e-ledger}). The hiding variant needs $\kappa \ge 14$, a stack of at least $2^{20}$ elements, for the reason below.

\paragraph{Zero knowledge.} The argument is for an honest verifier and statistical, then Fiat-Shamir in the random-oracle model. Every scalar sent before the lane fold's end is uniform, each under its own key. What the opening reveals besides is simulated in three parts.
\begin{itemize}[itemsep=1pt,topsep=2pt]
\item \emph{Level-$0$ answers.} For distinct $x_1, \dots, x_t \in \dom_0$ with $t \le \npad$, the map $\padc \mapsto \bigl(\Wn{\kappa}(x_i + \padshift) \sum_{j < \npad} \padc[j] \novel_j(x_i)\bigr)_i$ is onto $\K^t$: the $\novel_j$, $j < \npad$, span the polynomials of degree below $\npad$, which take any $t \le \npad$ values at $t$ points, and every factor is nonzero. So any $\npad$ symbols of a padded row are uniform whatever its message, the $(\nq_0 + 1)$-wise independence of the padded code~\cite{DHRR26,Dia25}, and a query message opens at most $\nq_0$ distinct columns.
\item \emph{Roots.} Given the opened columns, one more symbol of each row is still uniform ($\nq_0 + 1 \le \npad$), so each unopened leaf keeps one spare uniform symbol of $\K$ per row, $64$ bits of fresh entropy, and the root and the authentication paths are simulated in the random-oracle model with no salt.
\item \emph{After the lane fold.} No claim touches $\rlane$, so the lane fold gives $f^{(1)} = \sum_{u} \eq(\fc, u)\, f(u, \cdot)$ the term $\eq(\fc, \rlane)\, f(\rlane, \cdot)$, a nonzero multiple in $\E$ of a uniform vector of $\K^{2^\kappa}$ unless some $\fc_j$ is $0$ or $1$, which has probability at most $2\ell_0/|\E|$. Every later message is a linear functional of $f^{(1)}$, or for a root a random-oracle hash of the encoding of a fold of it: $c'$, the round and intro polynomials, the OOD answers, the values $c_q$, the opened rows of the levels past $0$, and $\ffinal$. An $\E$-valued functional is three $\K$-linear functionals of $f(\rlane, \cdot)$, and the configurations keep their number at most $2^{\kappa - 1}$, three per $\E$ scalar revealed, which is what fixes $\kappa \ge 14$; under that condition the random row makes them uniform, as VEIL's random column does~\cite[Lemma~4.7]{DHRR26}. The opened level-$0$ columns and $\padfold$ together determine the $c_q$, and reveal nothing more of the messages. This last step adapts VEIL's lemma to an $\E$-fold of $\K$-valued rows, with a random row of $\K$ and fold coefficients in $\E$: it is argued here, not proven.
\end{itemize}


%=====================================================================
\subsection{The Johnson bound}\label{app:johnson}
%=====================================================================

We prove Theorem~\ref{thm:johnson-interleaved} in its general form (arbitrary code).

\begin{theorem}[Johnson bound, agreement form]\label{thm:johnson-general}
Let $\mathcal{C} \subseteq \Sigma^\blen$ be a code, over any alphabet $\Sigma$, in which two distinct codewords agree on at most $\rate \blen$ coordinates, and let $\slack \in (0, 1 - \sqrt{\rate})$. Then every word $\mathbf{c} \in \Sigma^\blen$ has at most $\tfrac{1}{2\slack\sqrt{\rate}}$ codewords agreeing with it on at least $(\sqrt{\rate} + \slack)\, \blen$ coordinates.
\end{theorem}

\begin{proof}
Let $U_1, \dots, U_L$ be distinct codewords, each agreeing with $\mathbf{c}$ on at least $(\sqrt{\rate}+\slack) \blen$ coordinates. For each $j$, fix a set $\aset_j$ of exactly $a := \lceil (\sqrt{\rate}+\slack) \blen \rceil$ coordinates on which $U_j$ agrees with $\mathbf{c}$, and write $\varphi := a/\blen$, so $\sqrt{\rate} + \slack \le \varphi \le 1$. For a coordinate $i$, let $d_i := |\{ j : i \in \aset_j \}|$; then $\sum_i d_i = L a = L \varphi \blen$.

Wherever two codewords both agree with $\mathbf{c}$ they agree with each other, so $|\aset_j \cap \aset_{j'}| \le \rate \blen$ for $j \neq j'$, and counting ordered pairs,
\[
\sum_i d_i (d_i - 1) \;=\; \sum_{j \neq j'} |\aset_j \cap \aset_{j'}| \;\le\; L(L-1)\, \rate \blen .
\]
Cauchy--Schwarz on the vectors $(d_1, \dots, d_\blen)$ and $(1, \dots, 1)$ gives $\bigl(\sum_i d_i\bigr)^2 \le \blen \sum_i d_i^2$, i.e.\ $\sum_i d_i^2 \ge (L \varphi \blen)^2 / \blen = L^2 \varphi^2 \blen$; the left side above, which equals $\sum_i d_i^2 - \sum_i d_i$, is therefore at least $L^2 \varphi^2 \blen - L \varphi \blen$. Dividing by $L\blen$ and rearranging,
\[
L \varphi^2 - \varphi \;\le\; (L-1)\, \rate
\qquad\Longleftrightarrow\qquad
L\, (\varphi^2 - \rate) \;\le\; \varphi - \rate .
\]
Finally $\varphi^2 - \rate \ge (\sqrt{\rate}+\slack)^2 - \rate = 2\slack\sqrt{\rate} + \slack^2 > 2\slack\sqrt{\rate}$ while $\varphi - \rate \le 1$, so $L \le \tfrac{1}{2\slack\sqrt{\rate}}$.
\end{proof}
